\chapter{Практика}


\section{Настройка интерфейсов роутеров}

\begin{lstlisting}[language=C, caption={Пример настройки интерфейсов роутеров}]
#R1
configure

interface gigabitethernet 1/0/1
ip address 10.10.10.2/30
ip firewall disable
exit

interface gigabitethernet 1/0/3
ip address 192.168.1.1/24
ip firewall disable
exit

do commit
do confirm
 
#R2
configure

interface gigabitethernet 1/0/1
ip address 10.10.50.2/30
ip firewall disable
exit

interface gigabitethernet 1/0/3
ip address 192.168.1.1/24
ip firewall disable
exit

do commit
do confirm

#R2
configure

interface gigabitethernet 1/0/1
ip address 10.10.10.1/30
ip firewall disable
exit

interface gigabitethernet 1/0/3
ip address 10.10.50.1/30
ip firewall disable
exit

do commit
do confirm
\end{lstlisting}


\section{Настройка интерфейсов ПК}

\begin{lstlisting}[language=C, caption={Пример настройки интерфейсов ПК}]
# PC1
ip 192.168.1.100/24 192.168.1.1

# PC2
ip 192.168.1.10/24 192.168.1.1
\end{lstlisting}

\section{Настройка коммутаторов}

Интерфейсы коммутаторов необходимо перевести в режим switchport

\begin{lstlisting}[language=C, caption={Пример настройки коммутаторов}]
configure

interface gigabitethernet 1/0/1-24
mode switchport
exit

do commit
do confirm
\end{lstlisting}


\section{Настройка IP на Debian}

Чтобы настроить интерфейс, нужно внести изменения в файле /etc/netplan/interfaces

\begin{figure}[H] 
    \centering
    \includegraphics[width=0.75\textwidth]{debian.png}
    \caption{Настройка debian}
\end{figure}

После этого выполнить ifup ens3

\section{Настройка IP на Docker}

В этом моменте произошел баг: IP не настраивается (могу пингануть любое устройство из сети докера, но его пингануть не могу). По этой причине не могу получить
доступ к сайту на Debian

\begin{figure}[H] 
    \centering
    \includegraphics[width=0.75\textwidth]{error.png}
    \caption{Настройка docker}
\end{figure}


\section{Настройка source NAT на R1}

\begin{lstlisting}[language=C, caption={Пример настройки snat на R1}]
# переход к настройке source nat (будет меняться ip пакетов, которые выходят в публичную сеть (R1(gi 1/0/3) -> R1(gi 1/0/1)))
nat source
# создаем набор правил (имя любое)
ruleset SNAT
# выбираем интерфейс, к которому будут применяться правила (порт для выхода в публичную сеть)
to interface gi1/0/1
# задаем правило с номером 10
rule 10
# определяем пул адресов, которые будут подменяться
match source-address address-range 192.168.1.1-192.168.1.254
# включение трансляции
action source-nat interface
# активация
enable

exit
exit
exit

commit
do confirm
\end{lstlisting}


Теперь весь трафик, идущий с сети 192.168.1.0/24 будет поступать на порт 1/0/3 R1, а на выходе в публичную сеть с интерфейса 1/0/1 ip будет подменяться на ip интерфейса

\section{Настройка статического маршрута}

Чтобы роутер R1 пересылал трафик в публичную сеть, необходимо задать ему маршрут

\begin{lstlisting}[language=C, caption={Пример настройки статического маршрута на R1}]
ip route 0.0.0.0/0 10.10.10.1
\end{lstlisting}

Весь трафик будет пересылаться маршрутизатору R3

\section{Настройка R2}

Настройка аналогична R1, поскольку адреса подсетей совпадают. \\

Отличие будет в настройке маршрута: следующим hop'ом будет адрес 10.10.50.1


\section{Проверка работоспособности}

Если всё работает верно, то маршрутизатор R3 должен получать от маршрутизатора R2 подмененный адрес 


\begin{figure}[H] 
    \centering
    \includegraphics[width=0.75\textwidth]{result.png}
    \caption{Результат}
\end{figure}

R3 получает от R2 пакета с адресом 10.10.10.1, а не из сети 192.168.1.0/24. Это значит, что NAT работает


\endinput